% Generated draft: requires mathematical and source review.
\documentclass[11pt]{article}
\usepackage[margin=25mm]{geometry}
\usepackage{fontspec}
\setmainfont{Latin Modern Roman}
\usepackage{amsmath,amssymb,mathtools,unicode-math}
\setmathfont{Latin Modern Math}
\usepackage{enumitem,needspace,xurl,hyperref}
\setlength{\parindent}{0pt}
\setlength{\parskip}{5pt}
\setlength{\emergencystretch}{3em}
\newcommand{\N}{\mathbb{N}}
\newcommand{\Z}{\mathbb{Z}}
\newcommand{\Q}{\mathbb{Q}}
\newcommand{\R}{\mathbb{R}}
\newcommand{\C}{\mathbb{C}}
\newcommand{\F}{\mathbb{F}}
\newcommand{\A}{\mathbb{A}}
\newcommand{\PP}{\mathbb{P}}
\newcommand{\E}{\mathbb{E}}
\newcommand{\eps}{\varepsilon}
\newcommand{\dd}{\,\mathrm{d}}
\begin{document}
\section*{UnsolvedMath ID: 30000250}
\textbf{Source problem code:} catalogue entry
\subsection{Definitions and mathematical statement}
Retain \(F^a\) for the space and \(M\) for the Frobenius object occurring in the source. The conjectural assertion is that the Frobenius invariants \(M^{\varphi=1}\) form a local system of \(\Q_p\)-vector spaces on \(F^a\). Neither the space, the underlying coefficient sheaf, the construction of \(M\), nor the prime-dependent hypotheses are defined in the record. The superscript \(a\) is kept as part of the source notation and is not interpreted as an algebraic closure, an integer power or a named period-domain locus without evidence.

If \(M\) is realized on an intended site of \(F^a\) with a Frobenius operator compatible with restriction, the invariant sheaf would be described by
\[
 \mathcal L(U)=M^{\varphi=1}(U)
 =\{s\in M(U):\varphi(s)=s\}
 =\ker(\varphi-1)(U).
\]
This presumes the relevant identifications if Frobenius is semilinear or acts through a pullback; these cannot be read off from the notation alone. To obtain \(\Q_p\)-vector spaces, scalar multiplication by \(\Q_p\) must preserve the invariant sections, for example because Frobenius fixes the coefficient copy of \(\Q_p\). No exactness or surjectivity of \(\varphi-1\) is asserted.

In a locally constant finite-rank sheaf interpretation, the desired conclusion says that each point has a neighborhood in the chosen site on which \(\mathcal L\) is isomorphic to a constant sheaf with fiber \(\Q_p^r\), for some finite rank \(r\), locally constant on the base. A \(p\)-adic local system can instead be expressed in an appropriate lisse or continuous-representation framework; the topology and coefficient conventions matter, and ordinary discrete locally constant sheaves on an unspecified site cannot be silently substituted. The record gives no rank, no integral lattice or tower of finite-level sheaves, and no site on which local triviality is meant.

The quantifier concerns the particular family \(M\) constructed in the missing source context, not every sheaf equipped with an arbitrary endomorphism. A precise theorem requires that construction and its admissibility hypotheses. The literature triage checked 2026-08-21 reports related \(p\)-adic local-system constructions but no theorem found proving this exact assertion in the stated generality. Its open label and the verification field's report of source checking remain supplied metadata, not independently confirmed claims.
\subsection{Short English statement}
The conjecture asks whether the sections fixed by Frobenius in the intended family \(M\) form a \(\Q_p\)-local system on \(F^a\). The construction of the family, its scalar structure, the base space and the topology defining a local system are absent, so they cannot be replaced by an arbitrary Frobenius module or an assumed site.
\subsection{Sources}
\begin{itemize}
\item Local UnsolvedMath registry, permanent ID 30000250, source code catalogue entry, ``Local Systems from Frobenius-Invariant Sections,'' including original assertion, cleaned question and verification metadata.
\item \emph{Algebraische Zahlentheorie (2006)}, Oberwolfach Reports, pp.\ 1547--1588, DOI 10.4171/owr/2005/28: \url{https://doi.org/10.4171/owr/2005/28}. This is the sole supplied literature source, labelled ``original problem report.'' The citation year is 2006, while the DOI and proposed-year metadata use 2005.
\item Literature triage checked 2026-08-21; local record updated 2026-09-16. No external source has been visited or independently checked here.
\end{itemize}
\end{document}
